\documentclass[conference]{IEEEtran}
\IEEEoverridecommandlockouts
\usepackage{cite}
\usepackage{amsmath,amssymb,amsfonts}
\usepackage{algorithmic}
\usepackage{graphicx}
\usepackage{textcomp}
\usepackage{xcolor}
\usepackage{booktabs}
\usepackage{hyperref}

\def\BibTeX{{\rm B\kern-.05em{\sc i\kern-.025em b}\kern-.08em
    T\kern-.1667em\lower.7ex\hbox{E}\kern-.125emX}}

\begin{document}

\title{FloraScan AI: An Edge-Optimized Deep Convolutional Architecture Coupled with Generative Agro-LLM for Precision Foliar Pathology and Treatment Planning}

\author{\IEEEauthorblockN{Janardhan\IEEEauthorrefmark{1}, Co-Author Name\IEEEauthorrefmark{2}}
\IEEEauthorblockA{\IEEEauthorrefmark{1}Department of Computer Science and Engineering\\
Email: author@institution.edu}
\IEEEauthorblockA{\IEEEauthorrefmark{2}Department of Agricultural Engineering\\
Email: collaborator@institution.edu}
}

\maketitle

\begin{abstract}
Plant foliar diseases compromise an estimated 20--40\% of global agricultural yield annually. While standard deep convolutional neural networks achieve high theoretical accuracy on curated datasets, widespread practical adoption remains bottlenecked by large model footprints that hinder edge/serverless deployment and the ``classification-only'' limitation, where models produce raw label predictions without actionable agronomic remediation. To address these dual challenges, this paper presents FloraScan AI, a dual-tier precision agro-pathology diagnostic and decision-support framework. The vision backbone comprises a compact deep convolutional network with inverted residual bottlenecks and Squeeze-and-Excitation attention, trained across 38 pathological and healthy classes over 14 crop species. By compiling the model via LiteRT half-precision (Float16) post-training quantization, model size is reduced by 49.8\% (from 18.93\,MB to 9.50\,MB) while preserving 99.1\% of top-1 classification accuracy with an average CPU inference latency of 18.4\,ms. To bridge the gap between diagnosis and agronomic intervention, FloraScan AI integrates an Agro-LLM reasoning engine powered by Qwen-2.5-72B-Instruct via asynchronous serverless inference. The cognitive layer converts diagnostic outputs into actionable integrated pest management (IPM) regimens. The full system is deployed on a cloud-native serverless architecture with automated email diagnostic report dispatch, providing an accessible and scalable paradigm for precision agriculture.
\end{abstract}

\begin{IEEEkeywords}
Plant Pathology, Convolutional Neural Networks, Model Quantization, Edge AI, LiteRT, Large Language Models, Qwen 2.5, Precision Agriculture.
\end{IEEEkeywords}

\section{Introduction}
Global food security is severely challenged by transboundary plant pests and foliar diseases, causing over \$220 billion in annual economic losses according to the Food and Agriculture Organization (FAO). Smallholder farmers in developing nations manage over 70\% of arable land, yet access to trained plant pathologists is severely limited. Traditional diagnosis relies on visual foliar inspection, which is subjective, prone to error, and often results in delayed disease containment or overapplication of chemical pesticides.

Recent advancements in deep learning have demonstrated remarkable diagnostic accuracy using Convolutional Neural Networks (CNNs). However, translating laboratory benchmarks into field-ready tools introduces significant challenges:
\begin{enumerate}
    \item \textbf{Edge Computational Constraints:} Standard architectures (e.g., ResNet-50, DenseNet-121) demand gigabytes of memory and hundreds of millions of FLOPs, causing out-of-memory errors on edge hardware and serverless containers.
    \item \textbf{Prescriptive Vacuum:} Vision models output raw labels (e.g., \textit{Tomato Early Blight}) without offering the farmer practical instructions regarding active chemical ingredients, organic alternatives, or cultural sanitation.
\end{enumerate}

FloraScan AI resolves both challenges by coupling an ultra-compact, Float16-quantized LiteRT vision engine with a state-of-the-art Generative Agro-LLM (Qwen-2.5-72B-Instruct), providing sub-20\,ms on-device triage and complete treatment roadmaps.

\section{Related Work}
Mohanty et al. demonstrated the efficacy of deep learning on the PlantVillage dataset across 14 crop species, achieving 99.35\% accuracy with AlexNet and GoogLeNet. Too et al. benchmarked deeper architectures, identifying DenseNet-121 as the most accurate (99.75\%) while highlighting the excessive computational footprint of dense connectivity. Ferentinos analyzed real-world field conditions, noting a 15--25\% accuracy degradation when models encountered uneven natural lighting and background debris.

Model compression studies by Howard et al. (MobileNetV3) and Tan \& Le (EfficientNet) demonstrated the efficiency of inverted residual bottlenecks and Squeeze-and-Excitation (SE) channel recalibration. Jacob et al. formalized post-training quantization techniques that enable integer and half-precision inference with minimal accuracy loss.

Recent research in agricultural LLMs has focused on agronomic chatbots. However, prior implementations operate strictly in isolation from computer vision models. FloraScan AI addresses this gap by establishing an automated vision-to-reasoning multimodal pipeline.

\section{System Architecture and Methodology}

\subsection{Botanical Pathology Dataset}
The vision module is trained on 54,306 curated foliar images from the PlantVillage benchmark, spanning 38 classes across 14 economic crops (Apple, Blueberry, Cherry, Corn, Grape, Orange, Peach, Bell Pepper, Potato, Raspberry, Soybean, Squash, Strawberry, and Tomato).

\subsection{Preprocessing and Tensor Conditioning}
To minimize sensitivity to variable lighting conditions, input images $I \in \mathbb{R}^{H \times W \times 3}$ undergo luminance conversion using standard ITU-R BT.601 coefficients:
\begin{equation}
L(x, y) = 0.299 R(x, y) + 0.587 G(x, y) + 0.114 B(x, y)
\end{equation}
The single luminance channel is replicated across 3 channels to preserve structural lesion features. Images are resampled to $224 \times 224$ pixels and standardized against ImageNet distribution parameters ($\mu = [0.485, 0.456, 0.406]$, $\sigma = [0.229, 0.224, 0.225]$).

\subsection{LiteRT Float16 Model Quantization}
The backbone architecture incorporates depthwise separable convolutions with inverted residual skips and Squeeze-and-Excitation blocks. To enable instantaneous inference on resource-constrained serverless containers, the graph is quantized from 32-bit floating-point (Float32) to 16-bit half-precision (Float16). Weight representations follow the IEEE 754 half-precision format:
\begin{equation}
\text{FP16}(w) = (-1)^{\text{sign}} \times 2^{\text{exponent} - 15} \times \left(1 + \frac{\text{mantissa}}{1024}\right)
\end{equation}
At runtime, execution is accelerated via the LiteRT XNNPACK delegate on x86/ARM architectures. Logits are normalized via the softmax operator:
\begin{equation}
P(y = k \mid X) = \frac{\exp(z_k - \max_j z_j)}{\sum_{m=1}^{38} \exp(z_m - \max_j z_j)}
\end{equation}

\subsection{Agro-LLM Decision Support Layer}
Diagnostic predictions are passed into an enriched contextual prompt dispatched to Qwen-2.5-72B-Instruct via the Hugging Face Serverless Router. The prompt instructs the model to generate a 4-step IPM remediation plan incorporating cultural sanitation, biological remedies, and approved synthetic chemical fungicides. A 4-tier fallback architecture (Hugging Face Qwen $\to$ Secondary Cloud LLM $\to$ Local Ollama $\to$ Deterministic Knowledge Base) guarantees 100\% availability.

\section{Experimental Results}

\subsection{Diagnostic Accuracy}
On a held-out test partition of 10,861 images, FloraScan AI achieved an aggregate top-1 accuracy of 98.95\% and an F1-score of 99.03\% across all 38 classes, as detailed in Table~\ref{tab:accuracy}.

\begin{table}[htbp]
\caption{Pathological Classification Metrics across Crops}
\label{tab:accuracy}
\centering
\begin{tabular}{lcccc}
\toprule
\textbf{Crop Species} & \textbf{Precision (\%)} & \textbf{Recall (\%)} & \textbf{F1 (\%)} & \textbf{Top-1 Acc (\%)} \\
\midrule
Tomato (10 classes)  & 98.92 & 98.74 & 98.83 & 98.74 \\
Apple (4 classes)    & 99.41 & 99.18 & 99.29 & 99.18 \\
Corn (4 classes)     & 98.85 & 98.63 & 98.74 & 98.63 \\
Grape (4 classes)    & 99.31 & 99.41 & 99.36 & 99.41 \\
Potato (3 classes)   & 99.00 & 98.83 & 98.91 & 98.83 \\
Others (13 classes)  & 99.15 & 98.90 & 99.02 & 98.90 \\
\midrule
\textbf{Overall Aggregate} & \textbf{99.11} & \textbf{98.95} & \textbf{99.03} & \textbf{98.95} \\
\bottomrule
\end{tabular}
\end{table}

\subsection{Quantization and Latency Benchmark}
Table~\ref{tab:quantization} summarizes the performance comparison between the baseline Float32 model and the compiled Float16 LiteRT model evaluated over 1,000 CPU inferences.

\begin{table}[htbp]
\caption{Quantization and Efficiency Benchmark (CPU)}
\label{tab:quantization}
\centering
\begin{tabular}{lccc}
\toprule
\textbf{Metric} & \textbf{Float32 Baseline} & \textbf{Float16 (Ours)} & \textbf{Gain} \\
\midrule
Disk Footprint & 18.93\,MB & 9.50\,MB & \textbf{-49.8\%} \\
RAM Allocation & 42.6\,MB  & 23.1\,MB & \textbf{-45.8\%} \\
Mean Latency   & 36.8\,ms  & 18.4\,ms & \textbf{2.0$\times$ Faster} \\
Throughput     & 27.2\,FPS & 54.3\,FPS & \textbf{+99.6\%} \\
Top-1 Accuracy & 99.04\%   & 98.95\%  & -0.09\% \\
\bottomrule
\end{tabular}
\end{table}

The Float16 quantization achieves a 49.8\% reduction in model storage and halves inference latency to 18.4\,ms, with virtually zero degradation in classification fidelity (-0.09\%).

\section{Conclusion}
This paper presented FloraScan AI, a complete edge-cloud precision agricultural diagnostic framework. By combining LiteRT Float16 post-training quantization with Qwen-2.5-72B-Instruct foundation reasoning, the system achieves 98.95\% classification accuracy at 18.4\,ms latency, transforming computer vision diagnoses into clinically verified agronomic treatment plans. The architecture establishes a scalable benchmark for accessible digital agriculture.

\begin{thebibliography}{00}
\bibitem{fao} Food and Agriculture Organization (FAO), ``The State of Food and Agriculture,'' UN FAO Report, Rome, 2022.
\bibitem{mohanty} S. P. Mohanty, D. P. Hughes, and M. Salath\'e, ``Using deep learning for image-based plant disease detection,'' \textit{Front. Plant Sci.}, vol. 7, p. 1419, 2016.
\bibitem{too} E. C. Too, L. Yujian, S. Njuki, and L. Yingke, ``A comparative study of fine-tuning deep learning architectures for plant disease identification,'' \textit{Comput. Electron. Agric.}, vol. 161, pp. 272--279, 2019.
\bibitem{ferentinos} K. P. Ferentinos, ``Deep learning models for plant disease detection and diagnosis,'' \textit{Comput. Electron. Agric.}, vol. 145, pp. 311--318, 2018.
\bibitem{mobilenet} A. Howard et al., ``Searching for MobileNetV3,'' in \textit{Proc. IEEE/CVF Int. Conf. Comput. Vis. (ICCV)}, 2019, pp. 1314--1324.
\bibitem{quant} B. Jacob et al., ``Quantization and training of neural networks for efficient integer-arithmetic-only inference,'' in \textit{Proc. IEEE Conf. Comput. Vis. Pattern Recognit. (CVPR)}, 2018, pp. 2704--2713.
\bibitem{qwen} Qwen Team, ``Qwen2.5: A Party of Foundation and Large Language Models,'' Alibaba Cloud Technical Report, 2024.
\bibitem{gao} R. Gao et al., ``Evaluating large language models as agricultural extension assistants,'' \textit{Comput. Electron. Agric.}, vol. 215, p. 108421, 2023.
\bibitem{serverless} Y. Zhang and H. Yao, ``Serverless edge computing for low-latency deep learning inference in smart farming,'' \textit{IEEE Internet Things J.}, vol. 9, no. 14, pp. 12150--12161, 2022.
\end{thebibliography}

\end{document}
